\chapter{2009年湖南卷（文）}

\section{单选题}

\begin{problem}
\(\log_2\sqrt{2}\) 的值为（\quad）

\choices
    {\(-\sqrt{2}\)}
    {\(\sqrt{2}\)}
    {\(-\frac{1}{2}\)}
    {\(\frac{1}{2}\)}
\end{problem}

\begin{answer}
D.
\end{answer}

\begin{solution}
因为 \(\sqrt{2}=2^{\frac{1}{2}}\)，所以 \(\log_{2}\sqrt{2}=\log_{2}2^{\frac{1}{2}}=\frac{1}{2}\).
故选 D.
\end{solution}

\begin{problem}
抛物线 \( y^2 = -8x \) 的焦点坐标是（\quad）

\choices
    {\((2,0)\)}
    {\((-2, 0)\)}
    {\((4,0)\)}
    {\((-4,0)\)}
\end{problem}

\begin{answer}
B.
\end{answer}

\begin{solution}
抛物线的标准方程为 \(y^2=4px\)，与 \(y^2=-8x\) 对比得 \(4p=-8\)，即 \(p=-2\). 因而其焦点为 \((p,0)=(-2,0)\)，故选 B.
\end{solution}

\begin{problem}
设 \(S_{n}\) 是等差数列 \(\{a_{n}\}\) 的前 \(n\) 项和；已知 \(a_{2} = 3\)，\(a_{6} = 11\)，则 \(S_{7}\) 等于（\quad）

\choices
    {\(13\)}
    {\(35\)}
    {\(49\)}
    {\(63\)}
\end{problem}

\begin{answer}
C.
\end{answer}

\begin{solution}
设等差数列 \(\{a_n\}\) 的首项为 \(a_1\)，公差为 \(d\). 由
\(a_2=a_1+d=3\)，\(a_6=a_1+5d=11\)，相减得 \(4d=8\)，所以 \(d=2\)，进而 \(a_1=1\). 因此
\(a_7=a_1+6d=13\)，从而 \(S_7=\frac{7(a_1+a_7)}{2}=\frac{7(1+13)}{2}=49\).
故选 C.
\end{solution}

\begin{problem}
\(D\)，\(E\)，\(F\) 分别是 \(\triangle ABC\) 的边 \(AB\)，\(BC\)，\(CA\) 的中点，则（\quad）

\choices
    {\(\overrightarrow{AD} +\overrightarrow{BE} +\overrightarrow{CF} = 0\)}
    {\(\overrightarrow{BD} -\overrightarrow{CF} +\overrightarrow{DF} = 0\)}
    {\(\overrightarrow{AD} +\overrightarrow{CE} -\overrightarrow{CF} = 0\)}
    {\(\overrightarrow{BD} -\overrightarrow{BE} -\overrightarrow{FC} = 0\)}
\end{problem}

\begin{answer}
A.
\end{answer}

\begin{solution}
取点 \(A\)，\(B\)，\(C\) 的位置向量分别为 \(\bs{a}\)，\(\bs{b}\)，\(\bs{c}\). 由中点公式，\(\overrightarrow{AD}=\frac{\bs{b}-\bs{a}}{2}\)，\(\overrightarrow{BE}=\frac{\bs{c}-\bs{b}}{2}\)，\(\overrightarrow{CF}=\frac{\bs{a}-\bs{c}}{2}\).
于是 \(\overrightarrow{AD}+\overrightarrow{BE}+\overrightarrow{CF}=\frac{\bs{b}-\bs{a}+\bs{c}-\bs{b}+\bs{a}-\bs{c}}{2}=0\).
故选 A.
\end{solution}

\begin{problem}
某地政府召集 5 家企业的负责人开会，其中甲企业有 2 人到会，其余 4 家企业各有 1 人到会，会上有 3 人发言，则这 3 人来自 3 家不同企业的可能情况的种数为（\quad）

\choices
    {\(14\)}
    {\(16\)}
    {\(20\)}
    {\(48\)}
\end{problem}

\begin{answer}
B.
\end{answer}

\begin{solution}
若甲企业有 2 人发言，则这 3 人不可能来自 3 家不同企业，所以甲企业至多有 1 人发言.

当甲企业有 0 人发言时，从其余 4 家企业各选 1 人，共有 \(\mathrm C_4^3=4\) 种；当甲企业有 1 人发言时，先从甲企业的 2 人中选 1 人，再从其余 4 家企业中选 2 家，共有 \(2\mathrm C_4^2=12\) 种. 因而总数为 \(\mathrm C_4^3+2\mathrm C_4^2=4+12=16\).
故选 B.
\end{solution}

\begin{problem}
平面六面体 \(ABCD - A_{1}B_{1}C_{1}D_{1}\) 中，既与 \(AB\) 共面也与 \(CC_{1}\) 共面的棱的条数为（\quad）

\choices
    {\(3\)}
    {\(4\)}
    {\(5\)}
    {\(6\)}
\end{problem}

\begin{answer}
C.
\end{answer}

\begin{solution}
与棱 \(AB\) 共面的棱包括底面中与它相交或平行的棱，以及相应的上底面棱和侧棱，即
\(AB\)，\(BC\)，\(CD\)，\(DA\)，\(A_1B_1\)，\(C_1D_1\)，\(AA_1\)，\(BB_1\).

与棱 \(CC_1\) 共面的棱包括与它相交的棱、与它平行的棱以及上底面中相应的棱，即
\(CC_1\)，\(BC\)，\(CD\)，\(B_1C_1\)，\(C_1D_1\)，\(AA_1\)，\(BB_1\)，\(DD_1\).

两组棱的公共部分为 \(BC\)，\(CD\)，\(C_1D_1\)，\(AA_1\)，\(BB_1\)，共有 \(5\) 条. 故选 C.
\end{solution}

\begin{problem}
若函数 \( y = f(x) \) 的导函数在区间 \([a, b]\) 上是增函数，则函数 \( y = f(x) \) 在区间 \([a, b]\) 上的图象可能是 （\quad）

\choices
    {\choicebitmap[width=0.15\paperwidth]{img/2009/hunan_liberal/q7_optA.png}}
    {\choicebitmap[width=0.15\paperwidth]{img/2009/hunan_liberal/q7_optB.png}}
    {\choicebitmap[width=0.15\paperwidth]{img/2009/hunan_liberal/q7_optC.png}}
    {\choicebitmap[width=0.15\paperwidth]{img/2009/hunan_liberal/q7_optD.png}}
\end{problem}

\begin{answer}
A.
\end{answer}

\begin{solution}
导函数 \(f'(x)\) 在 \([a,b]\) 上递增，表示函数图象上各点的切线斜率随 \(x\) 增大而增大，所以图象应当向上弯曲，即函数在该区间上是凸函数. 四个图象中，只有第一幅图的切线斜率从左到右逐渐增大；第二幅图的斜率逐渐减小，第三幅图的斜率保持不变，第四幅图的斜率先增大后减小. 故选第一项.
\end{solution}

\begin{problem}
设函数 \( y = f(x) \) 在 \((- \infty, + \infty)\) 内有定义. 对于给定的正数 \(K\)，定义函数 \( f_{K}(x) = \begin{cases}
f(x), & f(x) \leqslant K,\\
K, & f(x) > K.
\end{cases}\) 取函数 \( f(x) = 2^{-|x|} \). 当 \( K = \frac{1}{2} \) 时，函数 \( f_{K}(x) \) 的单调递增区间为（\quad）

\choices
    {\((-\infty, 0)\)}
    {\((0, +\infty)\)}
    {\((-\infty, -1)\)}
    {\((1, +\infty)\)}
\end{problem}

\begin{answer}
C.
\end{answer}

\begin{solution}
当 \(x<0\) 时，\(f(x)=2^x\)；当 \(x\geqslant0\) 时，\(f(x)=2^{-x}\). 由 \(f(x)=\frac12\) 得 \(|x|=1\). 因而当 \(K=\frac12\) 时，
\[
f_K(x)=
\begin{cases}
2^x,&x\leqslant-1,\\
\frac12,&-1<x<1,\\
2^{-x},&x\geqslant1.
\end{cases}
\]
函数在 \((-\infty,-1)\) 上递增，在 \((-1,1)\) 上为常数，在 \((1,+\infty)\) 上递减. 题目给出的单调递增区间为 \((-\infty,-1)\)，故选 C.
\end{solution}

\section{填空题}

\begin{problem}
某班共 30 人，其中 15 人喜爱篮球运动，10 人喜爱乒乓球运动，8 人对这两项运动都不喜爱，则喜爱篮球运动但不喜爱乒乓球运动的人数为\fillinblank{}.
\end{problem}

\begin{answer}
\(12\).
\end{answer}

\begin{solution}
设喜爱篮球、乒乓球的学生集合分别为 \(A\)，\(B\). 由题意，\(|A|=15\)，\(|B|=10\)，且 \(|A\cup B|=30-8=22\). 因此 \(|A\cap B|=|A|+|B|-|A\cup B|=15+10-22=3\). 喜爱篮球但不喜爱乒乓球的人数为 \(|A|-|A\cap B|=15-3=12\).
\end{solution}

\begin{problem}
若 \(x > 0\)，则 \(x + \frac{2}{x}\) 的最小值为\fillinblank{}.
\end{problem}

\begin{answer}
\(2\sqrt{2}\).
\end{answer}

\begin{solution}
因为 \(x>0\)，所以 \(x+\frac{2}{x}-2\sqrt{2}=\left(\sqrt{x}-\frac{\sqrt{2}}{\sqrt{x}}\right)^2\geqslant0\).
等号成立当且仅当 \(\sqrt{x}=\frac{\sqrt{2}}{\sqrt{x}}\)，即 \(x=\sqrt{2}\). 因而最小值为 \(2\sqrt{2}\).
\end{solution}

\begin{problem}
在 \((1 + \sqrt{x})^4\) 的展开式中，\(x\) 的系数为\fillinblank{}. （用数字作答）
\end{problem}

\begin{answer}
\(6\).
\end{answer}

\begin{solution}
由二项式定理，\((1+\sqrt{x})^4=1+4\sqrt{x}+6x+4x\sqrt{x}+x^2\).
其中含 \(x\) 的项为 \(6x\)，所以 \(x\) 的系数为 \(6\).
\end{solution}

\begin{problem}
一个总体为 \(A\)，\(B\) 两层，用分层抽样方法从总体中抽取一个容量为 10 的样本. 已知 \(B\) 层中每个个体被抽到的概率都为 \(\frac{1}{12}\)，则总体中的个体数为\fillinblank{}.
\end{problem}

\begin{answer}
\(120\).
\end{answer}

\begin{solution}
设总体容量为 \(N\)，样本容量为 \(n=10\). 分层抽样按各层个体数所占比例分配样本，因此总体中每个个体被抽到的概率相同. 已知 \(B\) 层中每个个体被抽到的概率为 \(\frac{1}{12}\)，所以总体抽样比为 \(\frac{1}{12}\). 因而 \(\frac{10}{N}=\frac{1}{12}\)，
解得 \(N=120\).
\end{solution}

\begin{problem}
过双曲线 \(C: \frac{x^2}{a^2} - \frac{y^2}{b^2} = 1\)（\(a > 0\)，\(b > 0\)）的一个焦点作圆 \(x^2 + y^2 = a^2\) 的两条切线，切点分别为 \(A\)，\(B\). 若 \(\angle AOB = 120^\circ\) （ \(O\) 是坐标原点），则双曲线 \(C\) 的离心率为\fillinblank{}.
\end{problem}

\begin{answer}
\(2\).
\end{answer}

\begin{solution}
设双曲线的焦点为 \(F\)，焦距的一半为 \(c\)，则 \(F\) 在 \(x\) 轴上，且 \(OF=c\). 由于从 \(F\) 作圆的两条切线，故 \(OA\perp FA\)，从而 \(\triangle OAF\) 是直角三角形. 又 \(OA\)，\(OB\) 关于 \(x\) 轴对称，且 \(\angle AOB=120^\circ\)，所以 \(\angle AOF=60^\circ\). 因此 \(\cos60^\circ=\frac{OA}{OF}=\frac{a}{c}\)，即 \(c=2a\). 双曲线的离心率为 \(e=\frac{c}{a}=2\).
\end{solution}

\begin{problem}
在锐角 \(\triangle ABC\) 中，\(BC = 1\)，\(B = 2A\)，则 \(\frac{AC}{\cos A}\) 的值等于\fillinblank{}，\(AC\) 的取值范围为\fillinblank{}.
\end{problem}

\begin{answer}
\(2\)，\((\sqrt{2},\sqrt{3})\).
\end{answer}

\begin{solution}
设 \(A=\alpha\)，则 \(B=2\alpha\)，\(C=\pi-3\alpha\). 由于三角形为锐角三角形，须有 \(0<\alpha<\frac{\pi}{4}\)，且 \(0<\pi-3\alpha<\frac{\pi}{2}\)，
所以 \(\frac{\pi}{6}<\alpha<\frac{\pi}{4}\).

由正弦定理，\(BC\) 与 \(AC\) 分别为角 \(A\)，\(B\) 的对边，故 \(\frac{AC}{\sin2\alpha}=\frac{BC}{\sin\alpha}=\frac{1}{\sin\alpha}\). 于是 \(AC=\frac{\sin2\alpha}{\sin\alpha}=2\cos\alpha\)，从而 \(\frac{AC}{\cos A}=2\). 又因为 \(\frac{\pi}{6}<\alpha<\frac{\pi}{4}\)，且余弦函数在该区间上递减，所以 \(\sqrt{2}=2\cos\frac{\pi}{4}<AC<2\cos\frac{\pi}{6}=\sqrt{3}\).
因此 \(AC\) 的取值范围为 \((\sqrt{2},\sqrt{3})\).
\end{solution}

\begin{problem}
如图，两块斜边长相等的直角三角板拼在一起，若 \(\overrightarrow{AD} = x\overrightarrow{AB} + y\overrightarrow{AC}\)，则 \(x =\fillinblank{}\)，\(y =\fillinblank{}\).

\bitmapfigure[max width=0.5\linewidth]{img/2009/hunan_liberal/q15_fig1.png}
\end{problem}

\begin{answer}
\(1 + \frac{\sqrt{3}}{2}\)，\(\frac{\sqrt{3}}{2}\).
\end{answer}

\begin{solution}
设 \(AB=AC=s\). 由图中 \(\triangle ABC\) 在 \(A\) 处为直角且 \(\angle ACB=45^\circ\)，得 \(AB=AC=s\)，\(BC=\sqrt{2}s\). 另一块直角三角板为 \(\triangle EBD\)，其斜边为 \(ED\)，且两块三角板斜边相等，所以 \(ED=BC=\sqrt{2}s\). 因为 \(\angle EBD=90^\circ\)，\(\angle BED=60^\circ\)，故 \(BD=ED\sin60^\circ=\sqrt{2}s\cdot\frac{\sqrt{3}}{2}=\frac{\sqrt{6}}{2}s\).

以 \(A\) 为原点，取 \(\overrightarrow{AB}\)，\(\overrightarrow{AC}\) 所在直线分别为直角坐标系的两条坐标轴，则 \(B=(s,0)\)，\(C=(0,s)\). 由图可知 \(B\)，\(E\)，\(C\) 三点共线，且 \(\angle EBD=90^\circ\)，所以 \(BD\perp BC\). 而直线 \(BC\) 的斜率为 \(-1\)，所以直线 \(BD\) 的斜率为 \(1\). 结合图中 \(D\) 位于第一象限及 \(BD=\frac{\sqrt{6}}{2}s\)，得 \(\overrightarrow{BD}=\left(\frac{\sqrt{3}}{2}s,\frac{\sqrt{3}}{2}s\right)=\frac{\sqrt{3}}{2}\overrightarrow{AB}+\frac{\sqrt{3}}{2}\overrightarrow{AC}\). 因而 \(\overrightarrow{AD}=\overrightarrow{AB}+\overrightarrow{BD}=\left(1+\frac{\sqrt{3}}{2}\right)\overrightarrow{AB}+\frac{\sqrt{3}}{2}\overrightarrow{AC}\).
所以 \(x=1+\frac{\sqrt{3}}{2}\)，\(y=\frac{\sqrt{3}}{2}\).
\end{solution}

\section{解答题}

\begin{problem}
已知向量 \(\bs{a} = (\sin \theta, \cos \theta - 2\sin \theta)\)，\(\bs{b} = (1, 2)\).
\begin{enumerate}
\item 若 \(\bs{a} \parallel \bs{b}\)，求 \(\tan \theta\) 的值；
\item 若 \(|\bs{a}| = |\bs{b}|\)，\(0 < \theta < \pi\)，求 \(\theta\) 的值.
\end{enumerate}
\end{problem}

\begin{solution}
\begin{enumerate}
\item 由 \(\bs a\parallel\bs b\)，两向量的行列式为零，即
\[
\begin{vmatrix}
\sin\theta & \cos\theta-2\sin\theta\\
1 & 2
\end{vmatrix}=0
\]
所以 \(2\sin\theta-(\cos\theta-2\sin\theta)=0\)，即 \(4\sin\theta-\cos\theta=0\). 若 \(\cos\theta=0\)，上式不成立，故 \(\cos\theta\ne0\)，从而 \(\tan\theta=\frac{1}{4}\).

\item 由 \(|\bs a|=|\bs b|\)，得 \(\sin^2\theta+(\cos\theta-2\sin\theta)^2=1^2+2^2=5\)，整理得 \(\sin^2\theta-\sin\theta\cos\theta=1\).
当 \(\cos\theta=0\) 时，结合 \(0<\theta<\pi\)，可得 \(\theta=\frac{\pi}{2}\)，且它满足上式.

当 \(\cos\theta\ne0\) 时，两边同除以 \(\cos^2\theta\)，得 \(\tan^2\theta-\tan\theta=\sec^2\theta=1+\tan^2\theta\)，所以 \(\tan\theta=-1\). 在 \(0<\theta<\pi\) 内，故 \(\theta=\frac{3\pi}{4}\). 综上，\(\theta=\frac{\pi}{2}\) 或 \(\theta=\frac{3\pi}{4}\).
\end{enumerate}
\end{solution}

\begin{problem}
为拉动经济增长，某市决定新建一批重点工程，分别为基础设施工程，民生工程和产业建设工程三类. 这三类工程所含项目的个数分别占总数的 \(\frac{1}{2}\)，\(\frac{1}{3}\)，\(\frac{1}{6}\)，现在 3 名工人独立地从中任选一个项目参与建设. 求：
\begin{enumerate}
\item 他们选择的项目所属类别互不相同的概率；
\item 至少有 1 人选择的项目属于民生工程的概率.
\end{enumerate}
\end{problem}

\begin{solution}
设 3 名工人选择的项目类别分别为独立随机变量. 三类项目的选择概率依次为 \(\frac12\)，\(\frac13\)，\(\frac16\).
\begin{enumerate}
\item 3 人选择的类别互不相同，3 个类别的排列共有 \(3!\) 种，因此所求概率为 \(3!\times\frac12\times\frac13\times\frac16=\frac16\).

\item 先求没有人选择民生工程的概率. 每名工人选择非民生工程的概率为 \(1-\frac13=\frac23\). 由于 3 名工人独立选择，所求概率为 \(1-\left(\frac23\right)^3=1-\frac{8}{27}=\frac{19}{27}\).
\end{enumerate}
\end{solution}

\begin{problem}
如图，在正三棱柱 \(ABC - A_{1}B_{1}C_{1}\) 中，\(AB = 4\)，\(AA_{1} = \sqrt{7}\). 点 \(D\) 是 \(BC\) 的中点，点 \(E\) 在 \(AC\) 上，且 \(DE \perp A_{1}E\).
\begin{enumerate}
\item 证明：平面 \( A_{1}DE \perp \) 平面 \(ACC_{1}A_{1}\)；
\item 求直线 \(AD\) 和平面 \(A_{1}DE\) 所成角的正弦值.
\end{enumerate}
\bitmapfigure[max width=0.45\linewidth]{img/2009/hunan_liberal/q18_fig1.png}
\end{problem}

\begin{solution}
建立空间直角坐标系，取 \(A=(0,0,0)\)，\(B=(4,0,0)\)，\(C=(2,2\sqrt{3},0)\)，\(A_1=(0,0,\sqrt{7})\).
设 \(E\) 将 \(AC\) 按比例 \(\lambda\) 分点，则 \(E=(2\lambda,2\sqrt{3}\lambda,0)\).
由图可知 \(E\) 是 \(AC\) 的内点，故 \(0<\lambda<1\). 于是 \(\overrightarrow{DE}=(2\lambda-3,\sqrt{3}(2\lambda-1),0)\)，\(\overrightarrow{A_1E}=(2\lambda,2\sqrt{3}\lambda,-\sqrt{7})\).
由 \(DE\perp A_1E\)，有 \(\overrightarrow{DE}\cdot\overrightarrow{A_1E}=4\lambda(4\lambda-3)=0\).
所以 \(\lambda=\frac34\)，即 \(E=\left(\frac32,\frac{3\sqrt3}{2},0\right)\)，\(\overrightarrow{DE}=\left(-\frac32,\frac{\sqrt3}{2},0\right)\).

又 \(\overrightarrow{AC}=(2,2\sqrt3,0)\)，\(\overrightarrow{AA_1}=(0,0,\sqrt7)\)，
从而 \(\overrightarrow{DE}\cdot\overrightarrow{AC}=0\)，\(\overrightarrow{DE}\cdot\overrightarrow{AA_1}=0\).
直线 \(DE\) 在平面 \(A_1DE\) 内，且垂直于平面 \(ACC_1A_1\) 内的两条相交直线 \(AC\)，\(AA_1\)，故平面 \(A_1DE\) 垂直于平面 \(ACC_1A_1\).

平面 \(A_1DE\) 的一个法向量为
\[
\bs n=\overrightarrow{DE}\times\overrightarrow{A_1E}
=\left(-\frac{\sqrt{21}}2,-\frac{3\sqrt7}{2},-3\sqrt3\right)
\]
且 \(|\bs n|=4\sqrt3\). 又 \(\overrightarrow{AD}=(3,\sqrt3,0)\)，所以 \(\overrightarrow{AD}\cdot\bs n=-3\sqrt{21}\)，\(|\overrightarrow{AD}|=2\sqrt3\).
设直线 \(AD\) 与平面 \(A_1DE\) 所成的角为 \(\alpha\)，则
\[
\sin\alpha=\frac{|\overrightarrow{AD}\cdot\bs n|}{|\overrightarrow{AD}|\,|\bs n|}
=\frac{3\sqrt{21}}{2\sqrt3\cdot4\sqrt3}
=\frac{\sqrt{21}}8
\]
\end{solution}

\begin{problem}
已知函数 \( f(x) = x^3 + bx^2 + cx \) 的导函数的图象关于直线 \( x = 2 \) 对称.
\begin{enumerate}
\item 求 \( b \) 的值；
\item 若 \( f(x) \) 在 \( x = t \) 处取得极小值，记此极小值为 \( g(t) \)，求 \( g(t) \) 的定义域和值域.
\end{enumerate}
\end{problem}

\begin{solution}
由 \(f'(x)=3x^2+2bx+c\) 的图象关于直线 \(x=2\) 对称，得其对称轴满足 \(-\frac{2b}{2\cdot3}=-\frac{b}{3}=2\)，
所以 \(b=-6\).

此时 \(f'(x)=3x^2-12x+c=3(x-2)^2+c-12\). 若 \(c\geqslant12\)，则 \(f'(x)\geqslant0\)，函数 \(f(x)\) 单调递增，不在某点取得极小值. 若 \(c<12\)，则 \(f'(x)=0\) 的两个根为
\[
2-\sqrt{\frac{12-c}{3}}<2<2+\sqrt{\frac{12-c}{3}}
\]
且导函数在右侧根处由负变正，所以极小值点为
\[
t=2+\sqrt{\frac{12-c}{3}}>2
\]
反之，对任意 \(t>2\)，取 \(c=12t-3t^2<12\)，就有 \(f'(t)=0\)，且 \(t\) 是右侧根，故 \(f(x)\) 在 \(x=t\) 处取得极小值. 因此 \(g(t)\) 的定义域为 \((2,+\infty)\). 由 \(f'(t)=0\) 得 \(c=12t-3t^2\)，于是
\[
g(t)=f(t)=t^3-6t^2+ct
=t^3-6t^2+(12t-3t^2)t
=-2t^3+6t^2
\]
故 \(g(t)\) 的定义域为 \((2,+\infty)\). 对 \(t>2\)，\(g'(t)=-6t^2+12t=-6t(t-2)<0\). 函数 \(g(t)\) 在 \((2,+\infty)\) 上递减，且 \(\lim_{t\to2^+}g(t)=8\)，\(\lim_{t\to+\infty}g(t)=-\infty\).
因此值域为 \((-\infty,8)\).
\end{solution}

\begin{problem}
已知椭圆 \(C\) 的中心在原点，焦点在 \(x\) 轴上，以两个焦点和短轴的两个端点为顶点的四边形是一个面积为 8 的正方形 （记为 \(Q\)）.
\begin{enumerate}
\item 求椭圆 \(C\) 的方程；
\item 设点 \(P\) 是椭圆 \(C\) 的左准线与 \(x\) 轴的交点，过点 \(P\) 的直线 \(l\) 与椭圆 \(C\) 相交于 \(M\)，\(N\) 两点. 当线段 \(MN\) 的中点落在正方形 \(Q\) 内 （包括边界） 时，求直线 \(l\) 的斜率的取值范围.
\end{enumerate}
\end{problem}

\begin{solution}
设椭圆方程为 \(\frac{x^2}{a^2}+\frac{y^2}{b^2}=1\)，其中 \(a>b>0\)，
焦半距为 \(c\)，则 \(a^2=b^2+c^2\). 正方形 \(Q\) 的四个顶点为两个焦点和两个短轴端点，即 \((\pm c,0)\)，\((0,\pm b)\). 其对角线长分别为 \(2c\)，\(2b\)，所以 \(2bc=8\).
由于正方形两条对角线相等，得 \(b=c\)，从而 \(b=c=2\)，\(a^2=8\). 因此
\[
 C:\frac{x^2}{8}+\frac{y^2}{4}=1
\]

椭圆的左准线为 \(x=-\frac{a^2}{c}=-4\)，故 \(P=(-4,0)\). 设直线 \(l\) 的斜率为 \(k\)，则 \(y=k(x+4)\).
代入椭圆方程，得交点横坐标满足
\[
(1+2k^2)x^2+16k^2x+32k^2-8=0
\]
设两根为 \(x_1\)，\(x_2\)，由韦达定理，弦 \(MN\) 的中点 \(G=(x_0,y_0)\) 满足 \(x_0=\frac{x_1+x_2}{2}=-\frac{8k^2}{1+2k^2}\)，\(y_0=k(x_0+4)=\frac{4k}{1+2k^2}\).
两交点存在的条件为 \(\Delta=32(1-2k^2)>0\)，
即 \(|k|<\frac1{\sqrt2}\).

正方形 \(Q\) 的边界由 \(|x|+|y|=2\) 给出. 因为 \(x_0\leq0\)，点 \(G\) 落在 \(Q\) 内（包括边界）等价于 \(y_0\leq x_0+2\)，\(y_0\geq-x_0-2\). 代入 \(x_0\)，\(y_0\)，分别得到 \(2k^2+2k-1\leq0\)，\(2k^2-2k-1\leq0\).
两不等式的解集交为 \(-\frac{\sqrt3-1}{2}\leq k\leq\frac{\sqrt3-1}{2}\).
该区间包含于 \(|k|<\frac1{\sqrt2}\)，所以其中每个斜率都确实使直线与椭圆有两个交点. 故直线 \(l\) 的斜率取值范围为
\[
\left[-\frac{\sqrt3-1}{2},\frac{\sqrt3-1}{2}\right]
\]
\end{solution}

\begin{problem}
对于数列 \(\{u_n\}\)，若存在常数 \(M > 0\)，对任意的 \(n \in \N^*\)，恒有 \(|u_{n+1} - u_n| + |u_n - u_{n-1}| + \cdots + |u_2 - u_1| \leqslant M\)，则称数列 \(\{u_n\}\) 为 \(B\)-数列.
\begin{enumerate}
\item 首项为 1，公比为 \( -\frac{1}{2} \) 的等比数列是否为 \(B\)-数列？请说明理由；
\item 设 \(S_{n}\) 是数列 \(\{x_{n}\}\) 的前 \(n\) 项和. 给出下列两组论断：
A 组：
\begin{circlelist}
\item 数列 \(\{x_n\}\) 是 \(B\)-数列；
\item 数列 \(\{x_n\}\) 不是 \(B\)-数列；
\end{circlelist}
B 组：
{\renewcommand{\circled}[2][]{\tikz[baseline=(char.base)]{\node[shape=circle,draw,inner sep=0.8pt](char){\small \number\numexpr#2+2\relax};}}
\begin{circlelist}
\item 数列 \(\{S_n\}\) 是 \(B\)-数列；
\item 数列 \(\{S_n\}\) 不是 \(B\)-数列.
\end{circlelist}
}
请以其中一组中的一个论断为条件，另一组中的一个论断为结论组成一个命题. 判断所给命题的真假，并证明你的结论；
\item 若数列 \(\{a_n\}\) 是 \(B\)-数列，证明：数列 \(\{a_{n}^{2}\}\) 也是 \(B\)-数列.
\end{enumerate}
\end{problem}

\begin{solution}
\begin{enumerate}
\item 记该等比数列为 \(u_n=\left(-\frac12\right)^{n-1}\). 则 \(|u_{j+1}-u_j|=\left|\left(-\frac12\right)^j-\left(-\frac12\right)^{j-1}\right|=\frac32\left(\frac12\right)^{j-1}\).
因此，对任意 \(n\geq1\)，定义中的和为
\[
\sum_{j=1}^{n}|u_{j+1}-u_j|
=\frac32\sum_{j=0}^{n-1}\left(\frac12\right)^j
=3\left(1-2^{-n}\right)<3
\]
取 \(M=3\)，该数列是 \(B\)-数列.

\item 例如取命题“若数列 \(\{x_n\}\) 是 \(B\)-数列，则数列 \(\{S_n\}\) 是 \(B\)-数列”. 该命题为假：取 \(x_n=1\)，则 \(\{x_n\}\) 是 \(B\)-数列，而 \(S_n=n\)，有 \(\sum_{j=1}^{n}|S_{j+1}-S_j|=\sum_{j=1}^{n}1=n\)，
故 \(\{S_n\}\) 不是 \(B\)-数列，所以该命题为假.

另取命题“若数列 \(\{S_n\}\) 是 \(B\)-数列，则数列 \(\{x_n\}\) 是 \(B\)-数列”. 设 \(\{S_n\}\) 的定义中所用常数为 \(M\)，则
\[
|x_{n+1}|+|x_n|+\cdots+|x_2|
=|S_{n+1}-S_n|+|S_n-S_{n-1}|+\cdots+|S_2-S_1|\leq M
\]
又由 \(x_1=S_1\)，对任意 \(n\geq1\)，有
\[
\sum_{j=1}^{n}|x_{j+1}-x_j|
\leq |x_{n+1}|+2\sum_{j=2}^{n}|x_j|+|x_1|
\leq 2M+|x_1|
\]
所以 \(\{x_n\}\) 是 \(B\)-数列. 这给出了一个真的命题.

\item 设 \(\{a_n\}\) 为 \(B\)-数列，取常数 \(M>0\) 使 \(\sum_{j=1}^{n}|a_{j+1}-a_j|\leq M\). 于是对任意 \(n\)，\(|a_n|\leq |a_1|+\sum_{j=1}^{n-1}|a_{j+1}-a_j|\leq |a_1|+M\).
令 \(K=|a_1|+M\)，则 \(|a_n|\leq K\). 因此 \(|a_{j+1}^2-a_j^2|=|a_{j+1}+a_j|\,|a_{j+1}-a_j|\leq 2K|a_{j+1}-a_j|\).
对上式从 \(j=1\) 到 \(n\) 求和，得
\[
\sum_{j=1}^{n}|a_{j+1}^2-a_j^2|
\leq 2K\sum_{j=1}^{n}|a_{j+1}-a_j|\leq 2KM
\]
故 \(\{a_n^2\}\) 是 \(B\)-数列.
\end{enumerate}
\end{solution}
